\section{Regulatory Framework}

\label{sec:compliance}

\subsection{Isle of Man Money Transmitter License}

Lux operates the network's custodial and gateway functions under a
\textbf{Class 8 Money Transmitter License} issued by the Isle of Man
Financial Services Authority~\cite{iomfsa,cdax}. This class of
license authorizes bureau-de-change operations, payment services,
cheque cashing, and---critically for Lux---the issuance of
electronic money backed by real-world assets (Digital Safekeeping
Receipts, or D-SKRs).

The Isle of Man jurisdiction was selected for three reasons:

\begin{enumerate}[nosep]
  \item \textbf{Clear legal status for digital assets}: Manx law
        treats cryptocurrency as property for transfer purposes and
        money for transmitter licensing purposes, providing the
        clearest regulatory basis for token issuance.
  \item \textbf{EU and UK access via equivalence}: Isle of Man
        financial licenses enjoy mutual recognition with significant
        European and UK regulatory regimes.
  \item \textbf{Sophisticated regulator}: The IOMFSA has been issuing
        digital-money licenses since 2015 and has a track record of
        principled, technically competent supervision.
\end{enumerate}

\subsection{Identity Chain and DID}

The Identity Chain (I-Chain) stores Decentralized Identifier (DID)
records and verifiable credentials for every regulated participant
on the network. A participant's DID contains:

\begin{itemize}[nosep]
  \item A unique on-chain identifier (not derivable from the
        participant's personal data).
  \item Verifiable credentials issued by KYC providers, each with a
        cryptographic attestation and an expiry date.
  \item A revocation pointer that allows a KYC provider to revoke a
        credential on-chain if the participant's status changes.
  \item An accredited-investor flag (for US participants) or
        equivalent jurisdictional flags.
\end{itemize}

A smart contract that requires regulated transfers (e.g., a
security token) queries the I-Chain before permitting any transfer.
If the sender or recipient's DID does not carry a valid credential
for the relevant jurisdiction, the transfer reverts.

\subsection{Automated KYC/AML}

KYC/AML checks occur at three layers:

\begin{enumerate}[nosep]
  \item \textbf{Onboarding}: A new participant registers via a KYC
        provider (multiple providers are supported; the network is
        not beholden to any single vendor). The provider issues a
        verifiable credential to the participant's DID on the
        I-Chain.
  \item \textbf{Transaction}: Every regulated transfer queries the
        I-Chain to confirm that both parties hold valid credentials.
  \item \textbf{Ongoing}: KYC providers monitor their issued
        credentials for sanctions list changes, adverse media, and
        other ongoing screening events. A provider can revoke a
        credential at any time, immediately blocking the
        corresponding participant from further regulated transfers.
\end{enumerate}

%% ============================================================
